\section{Classification of Fields and Stress-Energy Tensor}

We now classify local fields according to their transformation properties under conformal mappings.

A field \(\Phi(z, \bar{z})\) is defined as a \textbf{conformal field} if it transforms homogeneously under independent global dilatations of the coordinates, namely \(z \to w = \lambda z\) and \(\bar{z} \to \bar{w} = \bar{\lambda} \bar{z}\). Under this scaling, the conformal field transforms as:
\[
    \Phi'(w, \bar{w}) = \lambda^{-h} \bar{\lambda}^{-\bar{h}} \Phi(z, \bar{z})
\]
The numbers \(h\) and \(\bar{h}\) are called the holomorphic and anti-holomorphic \textit{conformal weights} (sometimes also referred to as right and left conformal dimensions). Because the dilatation operator and the spin operator commute, any conformal field possesses both a definite \textbf{scale dimension} and a definite \textbf{spin}:
\[
    \Delta = h + \bar{h}, \quad s = h - \bar{h}.
\]
We can directly relate these quantum numbers to the generators of the global conformal algebra. The dilatation operator \(D\) and the rotation operator \(S\) are given by the sum and difference of the zero-modes of the Witt algebra:
\[
    D = \ell_0 + \bar{\ell}_0 \;, \quad S = \ell_0 - \bar{\ell}_0
\]
In terms of differential operators on the complex plane, these generators take the form \(D = z\partial_z + \bar{z}\partial_{\bar{z}}\) and \(S = z\partial_z - \bar{z}\partial_{\bar{z}}\). It is therefore straightforward to verify that applying these operators to a field yields precisely its scaling dimension and spin: \(D \to h + \bar{h} = \Delta\) and \(S \to h - \bar{h} = s\).
Consequently, under a pure global scaling \(z \to \lambda z\) and \(\bar{z} \to \lambda \bar{z}\), a field with well-defined weights transforms as:
\[
    \Phi(z, \bar{z}) \to \lambda^{-h} \lambda^{-\bar{h}} \Phi(z, \bar{z}) = \lambda^{-\Delta} \Phi(z, \bar{z})
\]
This explicitly shows that \(\Delta\) strictly governs the scaling behavior of the field.

A conformal field represents an element of the conformal basis within the algebra of local fields.

\subsubsection{Primary Fields}

Building upon this, a field \(\Phi(z, \bar{z})\) is classified as a \textbf{primary field} of conformal weights \((h, \bar{h})\) if it transforms covariantly under \textit{any} finite local conformal transformation \(z \mapsto w = f(z)\) and \(\bar{z} \mapsto \bar{w} = \bar{f}(\bar{z})\):
\[
    \Phi'(w, \bar{w}) = \left( \frac{\mathrm{d} w}{\mathrm{d} z} \right)^{-h} \left( \frac{\mathrm{d}\bar{w}}{\mathrm{d}\bar{z}} \right)^{-\bar{h}} \Phi(z, \bar{z}).
\]
This is the fundamental transformation law in two-dimensional CFT. We can easily derive its infinitesimal form. Consider a small local conformal transformation parameterized by
\[
    w = z + \epsilon(z), \quad \bar{w} = \bar{z} + \bar{\epsilon}(\bar{z}).
\]
The Jacobian factors expand as:
\[
    \frac{\mathrm{d} w}{\mathrm{d} z} = 1 + \partial_z \epsilon(z) \;, \quad \frac{\mathrm{d}\bar{w}}{\mathrm{d}\bar{z}} = 1 + \partial_{\bar{z}} \bar{\epsilon}(\bar{z}).
\]
Keeping only the first-order terms in \(\epsilon\), the conformal weight factors become:
\[
    \left( \frac{\mathrm{d} w}{\mathrm{d} z} \right)^{-h} \approx 1 - h \partial_z \epsilon \;, \quad \left( \frac{\mathrm{d}\bar{w}}{\mathrm{d}\bar{z}} \right)^{-\bar{h}} \approx 1 - \bar{h} \partial_{\bar{z}} \bar{\epsilon},
\]
and, at the same time, we can Taylor expand the field itself:
\[
    \Phi(z, \bar{z}) = \Phi(w - \epsilon, \bar{w} - \bar{\epsilon}) \approx \Phi(w, \bar{w}) - \epsilon \partial_w \Phi - \bar{\epsilon} \partial_{\bar{w}} \Phi.
\]
Finally, by substituting these expansions into the finite transformation law and isolating the variation \(\delta \Phi = \Phi'(z, \bar{z}) - \Phi(z, \bar{z})\), we obtain the infinitesimal transformation rule for a primary field:
\[
    \delta \Phi = -\epsilon(z) \partial_z \Phi - \bar{\epsilon}(\bar{z}) \partial_{\bar{z}} \Phi - h (\partial_z \epsilon(z)) \Phi - \bar{h} (\partial_{\bar{z}} \bar{\epsilon}(\bar{z})) \Phi.
\]
This transformation law shows that the field \(\Phi\) transforms as a tensor under all the local conformal transformations, with weights \(h\) and \(\bar{h}\). The primary fields are the subset of conformal fields that transform in this simple way, and they play a central role in the structure of the theory, as they are the building blocks from which all other fields can be generated by applying the generators of the conformal algebra.

\subsubsection{Secondary and Quasi-Primary Fields}

Many fields in a CFT, however, do not obey this strict local transformation rule. Fields that are not primary are collectively called \textbf{secondary fields}.

Among these, a special and very important subcategory exists: if a conformal field transforms covariantly \textit{only} under the restricted subset of global Möbius transformations (the \(\mathrm{SL}(2,\mathbb{C})\) group, where \(z \to w = \frac{az + b}{cz + d}\) with \(ad - bc = 1\)), it is called a \textbf{quasi-primary field}. Under such a transformation, a quasi-primary field behaves as:
\[
    \Phi'(w, \bar{w}) = (cz + d)^{2h} (\bar{c}\bar{z} + \bar{d})^{2\bar{h}} \Phi(z, \bar{z}).
\]
Every primary field is trivially quasi-primary, but the converse is generally false. A prominent example of a field that is quasi-primary but not primary is the energy-momentum tensor itself, which falls into the larger set of secondary fields because its transformation law acquires an anomalous shift under generic local mappings.

When we classify the operator content of a Conformal Field Theory, we will see that all fields naturally organize themselves into representations of the conformal algebra. The primary fields act as the fundamental building blocks of the theory. In representation theory, they correspond to the \textbf{highest-weight states}.

By acting on a primary field with the generators of local conformal transformations, we can generate an infinite family of quasi-primary and secondary fields known as the \textbf{descendants} of the primary field. In the algebraic language we will introduce later, a primary field is annihilated by the "lowering" generators (those with \(n > 0\)), while its descendants are constructed by repeatedly applying the "raising" operators (the generators with \(n < 0\)). Together, a primary field and all its generated descendants form what is called a complete conformal family or conformal module.

It is crucial to understand that the labels \((h, \bar{h})\) do not exclusively belong to primary fields. Even secondary or quasi-primary fields possess well-defined scaling properties locally. Under a general conformal map, the transformation law of any field always begins with the tensorial component \((\mathrm{d} w/\mathrm{d} z)^{-h}(\mathrm{d} \bar{w}/\mathrm{d}\bar{z})^{-\bar{h}}\Phi\), followed by anomalous terms (usually denoted by dots). The primary fields are simply the special building blocks for which these anomalous terms vanish completely.

Furthermore, the simple derivative of a primary field is generally not a primary field. If we were to take the derivative of the primary transformation law, the chain rule would generate second-order derivative terms (like \(\partial^2_z \epsilon\)), which spoil the strict definition of a primary field. However, taking derivatives shifts the conformal weights in a predictable way: the holomorphic derivative \(\partial_z \Phi\) of a primary field has weights \((h+1, \bar{h})\), while the anti-holomorphic derivative \(\partial_{\bar{z}} \Phi\) has weights \((h, \bar{h}+1)\).

\begin{example}[Transformation of the derivative of a field]
    Consider the holomorphic derivative \(\partial_z \Phi(z, \bar{z})\) of a primary field of weight \((h, \bar{h})\). Under a conformal map \(z \mapsto w(z)\), the derivative transforms according to the chain rule:
    \[
        \partial_z \Phi(z, \bar{z}) \mapsto \partial_z \left[ \left( \frac{dw}{dz} \right)^{-h} \Phi(w, \bar{w}) \right]
    \]
    Applying the product rule, we obtain:
    \[
        = \left( \frac{dw}{dz} \right)^{-h} \frac{dw}{dz} \partial_w \Phi(w, \bar{w}) - h \left( \frac{dw}{dz} \right)^{-h-1} \frac{d^2w}{dz^2} \Phi(w, \bar{w})
    \]
    The first term evaluates to \((dw/dz)^{-(h+1)} \partial_w \Phi\), which is exactly the transformation of a primary field with a shifted weight \((h+1, \bar{h})\). However, the second term introduces an anomaly proportional to the second derivative of the transformation, \(\partial^2_z w\). Because of this extra term, the derivative of a primary field is not a primary field itself, but rather a secondary field. \TODO{check computation.}
\end{example}

To conclude this geometric classification, it is useful to recall a few specific classes of conformal fields based on their weights:
\begin{itemize}
    \item \textbf{Scalar fields:} A scalar field has zero spin (\(s=0\)), which implies that its holomorphic and anti-holomorphic weights must be equal:
          \[
              h = \bar{h} = \frac{\Delta}{2}
          \]
    \item \textbf{Holomorphic (chiral) fields:} A field \(\Phi(z)\) that depends exclusively on \(z\) has \(\bar{h} = 0\). Its scale dimension and spin are completely determined by its holomorphic weight: \(\Delta = h\) and \(s = h\). Analytically, it satisfies the equation \(\partial_{\bar{z}} \Phi(z) = 0\).
    \item \textbf{Anti-holomorphic (anti-chiral) fields:} Conversely, a field \(\bar{\Phi}(\bar{z})\) depending only on \(\bar{z}\) has \(h = 0\), yielding \(\Delta = \bar{h}\) and \(s = -\bar{h}\). It satisfies the condition \(\partial_z \bar{\Phi}(\bar{z}) = 0\).
\end{itemize}
Physically, the differential equations satisfied by chiral and anti-chiral fields can be interpreted as continuity equations for a conserved current splitting into two helicity components, deeply tying these fields to underlying conservation laws (à la Noether) and fundamental symmetries.

\subsubsection{Stress-Energy Tensor}
We know that the stress-energy tensor \(T_{\mu\nu}\) is a local field, symmetric (or at least it can be made symmetric), and it is a conserved current, so it must satisfy the following conditions:
\[
    T_{\mu\nu}(x) = T_{\nu\mu}(x), \quad \partial^\mu T_{\mu\nu}(x) = 0, \quad T^\mu_\mu(x) = 0,
\]
where the continuity equation \(\partial^\mu T_{\mu\nu} = 0\) expresses the \textit{conservation of energy and momentum}, and the traceless condition \(T^\mu_\mu = 0\) is a consequence of \textit{scale invariance} at criticality (thus for Polyakov, since we have both rototranslation and scale invariance we know it to be conformal invariant).

We can again change to complex coordinates
\[
    \begin{dcases}
        z = x^1 + ix^2, \\
        \overline{z} = x^1 - ix^2,
    \end{dcases}
\]
and we can define the following components of the stress-energy tensor:
\[
    \begin{dcases}
        T(z,\overline{z}) = 2 \pi T_{zz}(z,\overline{z}) = 2 \pi( T_{11} - T_{22} + 2i T_{12} ), \\
        \overline{T}(z,\overline{z}) = 2 \pi T_{\overline{z}\overline{z}}(z,\overline{z}) = 2 \pi( T_{11} - T_{22} - 2i T_{12} ),
        \Theta(z,\overline{z}) = 2 \pi T_{z\overline{z}}(z,\overline{z}) = 2 \pi( T_{11} + T_{22} )=0.
    \end{dcases}
\]
where the \(2\pi\) factor comes from the original work of Polyakov, and it is just a convention. The last component \(\Theta\) is zero due to the traceless condition. The first two components \(T\) and \(\overline{T}\) are called the holomorphic and anti-holomorphic components of the stress-energy tensor, and they are the most important fields in the theory, since they generate the conformal transformations.

We can write the conservation equation \(\partial^\mu T_{\mu\nu} = 0\) in complex coordinates, and we find that it decouples into two independent equations:
\[
    \begin{dcases}
        \partial_{\overline{z}} T_{z z} + \partial_z T_{z\overline{z}} = 0, \\
        \partial_z \overline{T}(\overline{z}) + \dots = 0,
    \end{dcases}
\]
but since \(\Theta\) is zero, we can devolop
\[
    \dots
\]
and find out that \(T\) is undependent of \(\overline{z}\) and \(\overline{T}\) is independent of \(z\):
\[
    T=T(z), \quad \overline{T} = \overline{T}(\overline{z}), \quad \begin{dcases}
        \partial_{\overline{z}} T(z) = 0, \\
        \partial_z \overline{T}(\overline{z}) = 0.
    \end{dcases}
\]

Next let us take into account the action  as the integral of the lagrangian density, which is a local field, and thus we can write
\[
    S = \int \mathrm{d}^2 x \, \mathcal{L}(x) = \int \mathrm{d} z \mathrm{d} \overline{z} \, \mathcal{L}(z, \overline{z}),
\]
and since the action is supposed to have zero mass dimension, the lagrangian density must have scaling dimension \(\Delta = 2\). Thus, the lagrangian density is a local field with scaling dimension 2, and since we derive the stress-energy tensor from the lagrangian through
\[
    T = \frac{\partial \mathcal{L}}{\dots }\dots
\]
then we knwo the stress-energy tensor must have scaling dimension 2 as well, and thus it is a local field with scaling dimension 2. We would have expected it to have three components, but the traceless condition reduces it to two independent components, which are the holomorphic and anti-holomorphic components \(T(z)\) and \(\overline{T}(\overline{z})\).

We know that \(T\) is \((2,0)\) while \(\overline{T}\) is \((0,2)\), so they are not primary fields, since they do not transform tensorially under the conformal transformations, but they are secondary fields...

We know that a two point fucntion at the critical point should rescale by a power law, and thus we can write
\[
    \langle T(z)T(w)\rangle \sim \frac{1}{\vert z- w \vert^4 },
\]
thus if we do a mobius transformation when approaching the 4 oredr pole, it is fine, but a slightly more general transformation will give us some extra terms, which are not the transformation of a primary field, so \(T\) is not a primary field, but it is a secondary field. The same holds for \(\overline{T}\). \uline{Thus the stress-energy tensor is not a primary field.}

Now we are going to show two simple but fundamental results in two dimensional conformal field theory:
\begin{enumerate}
    \item The free boson in two dimensions as a conformal field theory.
    \item The free fermion in two dimensions as a conformal field theory.
\end{enumerate}











\section{Free Theories in Two Dimensions}

 [...]

\subsection{The Free Boson}

[...]
\[
    S = \frac{1}{8\pi} \int \mathrm{d}^2 x \, \partial_\mu \varphi(x) \partial^\mu \varphi(x),
\]
where the factor \(1/8\pi\) is just a convention, choosen by Ginsparg, but it is not important, the calculation are simpler.

    [comment on zero mass term]

A first result we can express with just field theory is
\[
    \langle \varphi(x)\varphi(y) \rangle = - \log \vert x - y \vert^2 + \text{const.},
\]
where \(x,y\) are two dimensional vectors, and the constant is an arbitrary constant which can be fixed by some normalization condition. By changing to complex coordinates, we can write
\[
    \langle \varphi(x)\varphi(y)\rangle = - \log (z - w ) + \log (\overline{z} - \overline{w}) + \text{const.},
\]
effectively decoupling the holomorphic and anti-holomorphic sectors. Furthermore if we consider the d'Alembert equation for the free boson
\[
    \box \varphi(z,\overline{z}) = \dots
\]
so that the more general solution of the equation of motion is given by the sum of a holomorphic and an anti-holomorphic, decoupled.

\begin{remark}
    The presence of the logarithm in the two-point function is a consequence of the fact that the free boson in two dimensions is a massless field, and thus it has a long-range correlation. This is not a conformal field, since we cannot put it in the set of fields of the algebra of local fields of the theory; actually Coleman in the Seventie wrote a paper showing that the massless boson in two dimensions is not a conformal field, but it does not mean that it can be a generator of a bigger algebra of fields, and thus it can be used to construct a conformal field theory, as we will see in the next chapters. If we insist with this lagrangian in thinking that \(\varphi\) as a foundamental field, we are simply wrong, but we can use this to build composite fields, which are conformal fields, and thus we can use the free boson to build a conformal field theory.
\end{remark}

Since the two components are totally decoupled, we can work with \(z\) and all the results will hold for \(\overline{z}\) as well.

We can define the following
\[
    \begin{aligned}
        \langle \partial \phi(z) \partial \phi(w) \rangle = \partial_z \partial_w \langle \phi(z) \phi(w) \rangle \\
        = - \partial_z \partial_w \log (z - w) = - \partial_w \frac{1}{z - w} = - \frac{1}{(z - w)^2},
    \end{aligned}
\]
and since there is a minus sign we can ``be russians about it'' and write the action with a minus sign in front, or instead define the following field
\[
    J(z) = i \partial \phi(z),
\]
for which we can write
\[
    \langle J(z) J(w) \rangle = \frac{1}{(z - w)^2},
\]
which has good properties, and it is a conformal field with scaling dimension \(\Delta = 1\) (since it transforms as the derivative of a primary field). It is a field that depends onaly on \(z\) and it is a \((1,0)\) field, so it is a chiral field. Its spin is \(s=1\) while for \(\overline{J}(\overline{z})\) we would find the same scaling dimension but spin \(s=-1\) and
\(\partial_\mu J^{\mu} = \partial J = \overline{\partial} \overline{J} =\overline{\partial} J = 0\), such that \(J\) and \(\overline{J}\) are conserved currents, and thus they are chiral currents. The charge implemented by the current do not necessary have a physical meaning, but we know that there have to be some symmetries in the theory.

Now let us try to study the operators product expansion of the field \(J(z)\) with itself, which is a fundamental tool in conformal field theory. We can write
\[
    \begin{aligned}
        J(z)J(w) = \contraction{}{J(z)}{}{J(w)} J(z) J(w) + \dots
    \end{aligned}
\]
where any part which is singular in this expansion is called a contraction (we are sort of generalizing the Wick theorem) and they will be represented by the connection among fields.
Thus up to regular terms, we can write
\[
    \langle J(z)J(w) \rangle = \frac{1}{(z-w)^2} = \contraction{}{J(z)}{}{J(w)} J(z) J(w),
\]
so that we can write the operator product expansion of \(J(z)\) with itself as
\[
    J(z)J(w) = \frac{1}{(z-w)^2} + \text{reg.}(z-w)^n,
\]
where the regular terms are less singular than the leading term, and they can be expanded in a power series in \(z-w\). Thus when we are computing the two point function of \(J(z)\) with itself, we are basically picking up the singular part of the operator product expansion, which is the contraction of the two fields. The regular part does not contribute to the two-point function, since it is regular when \(z\) approaches \(w\), and thus it does not have any singularity.\TODO{check.}
[...]
Then if we compute \(T(z)\) we can find
\[
    T(z) = - \frac{1}{2} : (\partial \phi(z))^2 : = \frac{1}{2} : J(z)^2 :,
\]
[...]
so we have to subtract the diverging part, which is basically the normal ordering of the fields (we are basically subtracting the vacuum expectation value of the product of fields, which is the singular part of the operator product expansion)...
    [...]

If we now compute the operator product expansion of \(T(z)\) with the conserved current \(J(w)\), we can find
\[
    T(z)J(w) = \frac{1}{2} : J(z)^2 : J(w),
\]
where if we apply the Wick theorem (contracting in all possible ways the fields) we can find\TODO{add normal ordering ::}
\[
    T(z)J(w) = \frac{1}{2} \contraction{}{J(z)}{}{J(z)} J(z) J(z) J(w) + \frac{1}{2} \contraction{}{J(z)}{}{J(w)} J(z) J(z) J(w),
\]
but since inside the normal ordering (the fields at most commute) we have only regular terms, the first term does not contribute, and we are left with
\[
    T(z)J(z) = \frac{1}{2} \contraction{}{J(z)}{}{J(w)} J(z) J(z) J(w) = \frac{J}{(z-w)^2} + \text{reg.}
\]

[\textit{OPE draw as z to w (points) }]

we know that the ope for two generic operators has to be of the form
\[
    \dots
\]
so we can expand \(J(z)\) in a Taylor series around \(w\) (since it is analytic in \(z\) and \(w\) is in the domain of analyticity)
\[
    J(z) = \dots
\]
and we can find
\[
    T(z)J(w) = \frac{J(w)}{(z-w)^2} + \frac{\partial J(w)}{z - w} + \text{reg.}
\]
which is a very important result, since it shows that \(J(w)\) is a primary field with scaling dimension \(\Delta = 1\), since the leading term in the OPE has a second order pole (and \(h=1\) is the coefficient in front of the leading term).

For the next OPE, we can compute \(T(z)T(w)\), which is the OPE of the stress-energy tensor with itself, and we can find\TODO{fix wick contractions}
\[
    \begin{aligned}
        T(z)T(w) & = \frac{1}{4} : J(z)J(z) : : J(w)J(w) :                                                                                                         \\
                 & = \frac{1}{4} \contraction{}{J(z)}{}{J(z)} J(z) J(z) : J(w)J(w) : + \frac{1}{4} \contraction{}{J(z)}{J(z)}{J(w)} J(z) J(z) : J(w)J(w) : + \dots
    \end{aligned}
\]
giving us terms (1), (2), and (3) which are the contractions of the fields, and we can compute them one by one:
\begin{enumerate}
    \item The first term gives us only regular terms, since the contraction of \(J(z)\) with itself is regular inside the normal ordering, so it does not contribute to the OPE.
    \item The second term gives us a contribution to the OPE, since the contraction of \(J(z)\) with \(J(w)\) gives us a singular term, and thus it contributes as
          \[
              \contraction{}{J(z)}{J(z)}{J(w)} J(z) J(z) : J(w)J(w) : = \frac{1}{(z-w)^2} : J(w)J(w) :,
          \]
          which can be expanded as \(J(z) = J(w) + (z-w) \partial J(w) + \dots\) in order to get
          \[
              \contraction{}{J(z)}{J(z)}{J(w)} J(z) J(z) : J(w)J(w) : = \frac{:J^2(w):}{(z-w)^2} + \frac{:J(w)\partial J(w):}{z-w} + \text{reg.}
          \]
          One can compute \(2 :J\partial J: = \partial :J^2: \) after some algebra, so we can write the contribution of this term as
          \[
              \contraction{}{J(z)}{J(z)}{J(w)} J(z) J(z) : J(w)J(w) : = \frac{2 :J^2(w):}{(z-w)^2} + \frac{\partial :J^2(w):}{z-w} + \text{reg.}
          \]
          but in the first term we can recognize the stress-energy tensor \(T(w) = \frac{1}{2} :J^2(w):\), and in the second term we can recognize the derivative of the stress-energy tensor \(\partial T(w) = \frac{1}{2} \partial :J^2(w):\), so we can write the contribution of this term as
          \[
              \contraction{}{J(z)}{J(z)}{J(w)} J(z) J(z) : J(w)J(w) : = \frac{2 T(w)}{(z-w)^2} + \frac{\partial T(w)}{z-w} + \text{reg.},
          \]
          and if the term (3) does not give us any contribution, we would have found that \(T(z)\) is a primary field with scaling dimension \(\Delta = 2\), but we will see that this is not the case.
    \item [...]

\end{enumerate}
\[
    T(z)T(w) = \frac{1}{2} \frac{1}{\vert z-w \vert^2} + \frac{2 T(w)}{(z-w)^2} + \frac{\partial T(w)}{z-w} + \text{reg.},
\]
so we have an extra term in the OPE, which is not the transformation of a primary field, and thus \(T(z)\) is not a primary field, but it is a secondary field.
In general this ope in general has an arbitrary constant \(c\) in front of the first term, which is called the \textbf{central charge} of the theory, and it is a very important parameter which characterizes the conformal field theory. In our case we have \(c=1\), but in general it can be any real number, and it can be negative as well.

We do not know yet the physical meaning of the central charge, but we will see that it is related to the number of degrees of freedom of the theory, and it is also related to the anomaly of the theory when we try to quantize it. It is caracteristic of the theory, since it comes from the OPE of the stress-energy tensor with itself, which is a fundamental field in the theory (coming directly from the lagrangian of the theory), and thus it is a parameter which characterizes the theory.

The classification of conformal field theories is the idea to find numbers such as the central charge which can determin univocally the theory, and thus to classify all the possible conformal field theories. This is not possible in general, but for example the \textbf{minimal models} (\(0<c<1\)) are a class of conformal field theories which are completely classified by the central charge and the scaling dimensions of the primary fields, and also completely solved, since we can compute all the correlation functions of the theory. The ising model...

\subsection{The Free Fermion}

[...]